\chapter{Del sonido al instrumento}\label{cap:instrumento}

El resultado de Fourierizer es un sonido. Para tocarlo hace falta un \emph{sampler}: un programa que,
cada vez que se presiona una tecla, reproduce ese sonido transpuesto a la nota correcta, con una
envolvente, y que puede hacerlo con muchas notas a la vez. Este capítulo describe el sampler de
Fourierizer y las reglas del \emph{tiempo real} que lo condicionan.

\section{Afinación: de la tecla a la transposición}

Una tecla MIDI $n$ corresponde a la frecuencia $440\cdot2^{(n-69)/12}\Hz$. El resultado suena tal
cual en la nota \emph{Root} $n_0$. Tocar la nota $n$ pide transponerlo
\[
  s=n-n_0\ \text{semitonos},\qquad r=2^{s/12}.
\]
Con \emph{Tape}, cada voz lee el resultado con paso $r\cdot\Fs^{\text{buf}}/\Fs^{\text{host}}$, que
también convierte la frecuencia de muestreo del resultado ($48$ kHz) a la del programa anfitrión.
Con \emph{Spectral}, cada voz lee a paso $\Fs^{\text{buf}}/\Fs^{\text{host}}$ una copia del resultado
ya transpuesta $s$ semitonos (\cref{cap:transformaciones}).

\paragraph{Interpolación.} Leer en posiciones no enteras exige interpolar. En tiempo real, el
sampler usa la interpolación cúbica de Hermite con cuatro muestras: continua, con derivada
continua y barata. A diferencia del remuestreador sinc del \cref{cap:transformaciones}, \emph{no}
filtra antes de acelerar. Al transponer hacia arriba ($r>1$), el contenido por encima de
$\Fs/(2r)$ se refleja por debajo de Nyquist (aliasing). En sonidos oscuros no se nota; en sonidos
brillantes transpuestos más de una octava, sí. La solución clásica son las \emph{mipmaps}: copias
prefiltradas del sonido por octava, entre las que se elige según la nota. Quedan como mejora
pendiente.

\section{La envolvente}

Cada voz multiplica su sonido por una envolvente $e(t)\in[0,1]$ de cuatro etapas (ADSR):
\[
  e(t)=\begin{cases}
    t/T_a, & 0\le t<T_a \quad\text{(ataque lineal)},\\[2pt]
    s+(1-s)\,\e^{-(t-T_a)/\tau}, & t\ge T_a,\ \tau=T_d/3 \quad\text{(decaimiento exponencial)},
  \end{cases}
\]
y al soltar la tecla, en el instante $t_r$, baja linealmente desde el nivel que tenía, $e(t_r)$, hasta
$0$ en $T_r$ segundos (el \emph{release}). El decaimiento es exponencial porque así decaen los
sistemas físicos (una cuerda, una membrana): pierden una fracción fija de energía por unidad de
tiempo, y el oído lo percibe como una caída ``pareja'' en decibeles. La constante $\tau=T_d/3$ hace
que en $T_d$ segundos la distancia al sostenido baje a $\e^{-3}\approx0.05$, el 5\,\%: el decaimiento
``llega'' al sostenido en el tiempo que marca la perilla.

Las perillas no son lineales. El ataque, por ejemplo, es $T_a=0.001+1.999\,v^2$ segundos para la
posición $v\in[0,1]$. La mitad de su recorrido cubre de $1$ a $500$ ms, la zona donde el oído
distingue más diferencias, y la otra mitad llega hasta $2$ s.

\begin{estudio}
\textbf{Envolventes típicas.}
\begin{itemize}
  \item \emph{Pads y texturas:} ataque largo (300 ms o más), sostenido alto, release largo. El sonido
        aparece y desaparece sin cortes.
  \item \emph{Plucks y percusivos:} ataque mínimo, decaimiento corto y sostenido en cero. El sonido
        dura lo que su decaimiento, sin importar cuánto se mantenga la tecla.
  \item Si un sonido Fourierizado ya trae su propia forma (por ejemplo, fades en la selección),
        usa ataque y release cortos para no tapar la que ya tiene.
\end{itemize}
\end{estudio}

\section{Polifonía}

El sampler tiene 24 voces. Cada nota ocupa una voz libre. Si no hay ninguna, se \emph{roba} la más
antigua: es la que lleva más tiempo sonando y suele ser la menos audible. El pedal de sostenido
(CC64) mantiene las notas soltadas hasta que se levanta. Las notas de la interfaz (el teclado de la
computadora) y las del programa anfitrión (MIDI) entran por caminos distintos, pero llegan a las
mismas voces.

\section{Las reglas del tiempo real}

El anfitrión pide audio en bloques de $B$ muestras: típicamente $B=256$ o $512$, es decir, $5.3$ o
$10.7$ ms a $48$ kHz. Cada bloque tiene que estar listo \emph{antes} de que se acabe el anterior. Si
no, la tarjeta de sonido se queda sin datos y se oye un chasquido. No basta con que el cálculo sea
rápido \emph{en promedio}: tiene que ser rápido \emph{siempre}. Por eso el hilo de audio sigue tres
reglas:
\begin{enumerate}
  \item \textbf{No pide ni libera memoria.} El administrador de memoria del sistema puede tardar un
        tiempo impredecible.
  \item \textbf{No espera a otros hilos.} Si toma un candado que otro hilo tiene, queda a merced de
        ese hilo, que puede tener menos prioridad (la \emph{inversión de prioridad}).
  \item \textbf{No lee ni escribe archivos,} por la misma razón.
\end{enumerate}

Las consecuencias en el diseño:
\begin{itemize}
  \item \emph{Un resultado nuevo} llega por un intercambio de punteros protegido con un candado que
        el hilo de audio solo \emph{intenta} tomar. Si está ocupado, lo toma en el bloque
        siguiente.
  \item \emph{Los buffers viejos} nunca se liberan en el hilo de audio. Un ``depósito de
        liberación'' conserva una referencia a cada buffer, y otro hilo lo libera cuando nadie más lo
        usa.
  \item \emph{Las notas de la interfaz} llegan por una cola circular sin candados.
\end{itemize}

\section{El banco de copias de Spectral}

Un phase vocoder de calidad es demasiado caro para correr por nota en el hilo de audio. Como el
resultado es fijo, Fourierizer calcula las copias transpuestas \emph{antes}: una por semitono, en dos
hilos de fondo. Primero van las notas que tocas y después la octava alrededor del Root. Hay un
presupuesto de memoria de 1 GB, y si se llena se descarta la copia menos usada recientemente.

Cada copia se produce por pedazos: un primer pedazo de 4096 muestras y luego de 16384. Una voz
puede empezar a tocar una copia que todavía se está calculando: solo tiene que no rebasar la
última muestra lista. Como el cálculo corre entre 8 y 40 veces más rápido que el tiempo real, la
reproducción nunca alcanza al cálculo. Una nota nueva espera, cuando mucho, su primer pedazo:
alrededor de un bloque.

Al \emph{exportar} (render fuera de línea), el anfitrión pide bloques más rápido que el tiempo real
y no importa esperar. En ese modo, la voz sí espera a que su copia esté lista, y el archivo
exportado sale exacto, sin depender de qué tan rápida sea la computadora.

\section{Estado y parámetros}

Todo lo que define el instrumento se guarda en el proyecto del anfitrión: los sonidos (en FLAC
24-bit, cada audio distinto una sola vez), las selecciones, las tomas y la petición del último
render. Al abrir el proyecto, el plugin vuelve a calcular el resultado en segundo plano, así que
suena aunque nadie abra su ventana. Los controles del sampler (ADSR, Root, volumen, Tape o Spectral
y la ventana de Spectral) son \emph{parámetros del anfitrión}: se pueden automatizar y ligar a un
controlador MIDI.

\begin{estudio}
\textbf{Elegir el Root.} Es la nota en la que el resultado suena intacto. Conviene que sea la nota real
del sonido, si la tiene (la vista de análisis del resultado la muestra), y que quede en el centro del
rango que vas a tocar. Así las transposiciones son pequeñas en ambas direcciones, que es donde Tape
y Spectral suenan mejor.
\end{estudio}
